The \peripheral{DAC} peripheral consists of two independent 14-bit R-2R DACs with a DMA interface for interrupt-driven data output. Its main features are listed below:

\begin{itemize}
	\item Sampling rate up to 700 kSPS
	\item Voltage range from 0~\si{\volt} to 2.39~\si{\volt}, enforced by bandgap generator
	\item 14-bit resolution, $V_{LSB}$ = 146~\si{\micro\volt}
	\item Output resistance of $\approx200$~\si{\ohm}
	\item Single conversion mode
	\item DMA to automatically fetch new samples directly from memory, with configurable buffer size
	\item Interrupt indicating the DMA buffer is empty
	\item DAC value can be left-aligned or right-aligned
\end{itemize}



\subsection{DAC Pins}

The \peripheral{DAC} peripheral has to pins, \pin{DAC0} and \pin{DAC1}, which are the two independent DAC outputs. When the DAC is enabled with \bitfield{DACxEN}, the pin is a low impedance ($\approx200$~\si{\ohm}) driven output. When the DAC is disabled, the output is pulled to ground with a large pulldown impedance of 400~\si{\kilo\ohm}. WARNING: Do NOT apply a voltage outside the bounds of 0.0 V to 2.5 V to any analog pin! Doing so can permanently damage the internal circuits.



\subsection{R-2R DAC Architecture}

The R-2R DAC architecutre is so named for its repetitive chain of two resistors, $R$ and $2R$, for every added bit of resolution (shown in Figure \ref{fig:r2r-dac-architecture}). Each bit $b_i$ in the DAC input value has an associated voltage $v_{bi}$, which is $V_{ref}$ when the bit $b_i$ is 1 or grounded when the bit is 0. Absent an in-depth circuit analysis, the result is that each bit adds voltage $2V_{ref}/2^{N-i}$ to the output, which is:

\begin{equation*}
	V_{DAC} = \frac{2V_{ref}}{2^N} \sum_{i=0}^{N-1} b_i 2^i
\end{equation*}
where $N$ is the number of bits in the DAC value, and $i$ is the bit number with 0 being the LSB and $N - 1$ being the MSB.

\begin{figure}[htpb]
	\centering
	\includesvg{figures/r2r-dac-architecture}
	\caption{R-2R DAC Architecture}
	\label{fig:r2r-dac-architecture}
\end{figure}

It should be noted that the number of bits $N$ cannot be increased indefinitely due to several physical limitations. First, each resistor introduces a thermal noise power spectral density of $V^2_R(f) = 4kTR$, where $R$ is the resistance, $k$ is the Boltzmann constant, and $T$ is absolute temperature. Thus, for some choice of $R$, there is a threshold for $N$ past which $V_{LSB}$ will be smaller than the total noise contribution from the resistors, which renders any additional bits overpowered by noise and effectively useless. Furthermore, the R-2R ladder assumes that each resistance is exactly equal to $R$ or $2R$, which does not reflect the reality of the manufacturing process. Due to small variations during fabrication, the resistors in the R-2R ladder will exhibit a statistical distribution of resistances. For some value of $N$ large enough, these differences can lead to output voltage swings greater or less than $2V_{ref}/2^{N-i}$ for each bit $b_i$, which again renders additional bits useless.

For the \peripheral{DAC} peripheral, the number of bits $N$ is 14, and the reference voltage $V_{ref}$ is generated by a bandgap reference circuit at 1.196~\si{\volt}. The output opamp attached to the end of the R-2R ladder is configured in feedback to double the voltage seen at its positive input. Therefore, the maximum voltage the DAC can output is $2 V_{ref}$, and the minimum is 0~\si{\volt}.



\subsection{DAC Analog Bias Setup}

The global bias generator must be enabled and calibrated to an appropriate bias level before the DAC can be used. The \bitfield{BIASEN} bit enables the global bias generator, and the \bitfield{BIASADJ} bit field controls the bias level. To determine the appropriate value of \bitfield{BIASADJ}, the user should systematically characterize the analog circuits for all possible calibration vectors and pick the vector that best suits the application. (A nominal value of \bitfield{BIASADJ} = 37 is suggested from simulation data.) Then, the bandgap generator must be enabled with the \bitfield{BGEN}. All three of these bitfields are contained in the \register{BIASCR} register.

The desired DAC(s) must also be enabled with the \bitfield{DACxEN} bits in the \register{DACxCR} register.



\subsection{Single Conversion Mode}

The simplest DAC mode of operation is single conversion mode, where the DAC produces a single output voltage whenever the \register{DACxVAL} register is written by the user. That voltage remains steady until the next time the \register{DACxVAL} register is written to. The user may change the value in \register{DACxVAL} as often as desired, but if this is done at a rate greater than 700~kSPS, then the DAC will not be able to settle on the current value before going to the next value.



\subsection{DAC Output Value Alignment}

The 14-bit DAC input value may be aligned either to the right (LSB) side or left (MSB) side of the 16-bit \register{DACxVAL} register and the DMA double buffer sample with the \bitfield{DACDAL} bit. When right-aligned, the DAC value will occupy the 14 least significant bits, and the 2 most significant bits will be \texttt{0}s. Likewise, when left-aligned, the DAC value will occupy the 14 most significant bits, and the 2 least significant bits will be \texttt{0}s. This is shown in Figure \ref{fig:dac-data-alignment}.

\begin{figure}[htpb]
	\centering
	\begin{subfigure}[b]{\linewidth}
		\begin{bytefield}[endianness=big, bitwidth=0.0625\linewidth]{16}
			\bitheader[lsb=0]{0-15} \\
			\bitbox{1}{\texttt{0}} & \bitbox{1}{\texttt{0}} & \bitbox{14}{14-bit DAC output value} \\
		\end{bytefield}
		\caption{Right-aligned (DAC value in LSB)}
		\label{fig:dac-data-alignment-right}
	\end{subfigure}
	\vspace*{0.25in}

	\begin{subfigure}[b]{\linewidth}
		\begin{bytefield}[endianness=big, bitwidth=0.0625\linewidth]{16}
			\bitheader[lsb=0]{0-15} \\
			\bitbox{14}{14-bit DAC output value} & \bitbox{1}{\texttt{0}} & \bitbox{1}{\texttt{0}} \\
		\end{bytefield}
		\caption{Left-aligned (DAC value in MSB)}
		\label{fig:dac-data-alignment-left}
	\end{subfigure}
	\caption{DAC Data Alignment}
	\label{fig:dac-data-alignment}
\end{figure}

Data alignment is desirable when the user wishes to place samples from the ADC directly into the DAC value. As the DAC value is 14 bits wide, placing a right-aligned ADC value into a right-aligned DAC value will cause the 4 most significant bits of the DAC to be \texttt{0}s, which may be undesireable. However, if both are left-aligned, then the MSB of the ADC value will align with the MSB of the DAC value, and only the 4 least significant bits of the DAC will be 0s.



\subsection{Periodic Sampling Mode with DMA}

The DAC has an inbuilt timer that allows it to read data directly from RAM to periodically update the DAC input value using a direct memory access (DMA) unit. The DAC clock determines the DAC sampling frequency, and cannot exceed about 700~kSPS without the internal DAC circuitry introducing the dominant pole in its frequency response. First, the DAC clock souce must be selected using the \bitfield{DACCS} bitfield from among the following options: MCLK, SMCLK, HFXT, or DCO0. Next, the DAC clock source frequency must be divided down using the \bitfield{DACCDIV} bitfield, which can divide the DAC clock source by powers of two from 1/1 to 1/128. Then, divided clock is further divided by 16, producing the DAC clock. The \register{DACxFS} register sets the number of DAC clock cycles are in between each DAC sample, yielding a sampling frequency of: 

\begin{equation*}
	F_{s} = \frac{f_{src}}{16 \times 2^\bitfield{DACCDIV} \times (1 + \register{DACxFS})}
\end{equation*}

The DMA reads each new DAC sample from RAM using a double buffer (or ping pong buffer), a configuration which allows both the DAC and CPU to access the samples in memory without memory collisions or errors. The double buffer is split into two halves: the first half, and the second half. At any given time while the DMA is activated, the DAC will be reading one half of the buffer and placing the samples contained there into the DAC value one by one, and the CPU will be allowed to write the next the samples to the other half. The CPU is not allowed to access the data in the half that the DAC is currently reading samples from. When the DAC empties its half of the buffer, the DMA sets the DAC buffer half empty interrupt flag \bitfield{DACBEIF} and then seizes control of the other half of the double buffer, allowing the CPU to write new samples to it. The DMA continues alternating buffers indefinitely unless the DMA is disabled. The user is responsible writing new samples in the other half of the buffer before the DAC begins reading that half.

Both DACs (DAC0 and DAC1) share the same DMA unit, meaning that the DMA updates their values simultaneously. Therefore, if both DACs are in DMA mode, they must both share the same sampling frequency and phase. Each DAC can independently enter DMA mode using the \bitfield{DACxDEN} bits, whereupon the DMA will immediately begin reading samples from RAM and pushing them to the selected DACs. (Note: the DACs must also be enabled with \bitfield{DACxEN}.) 

The first half of the DAC double buffer begins at address \texttt{0x11800}, and the second half begins at address \texttt{0x11C00}. Both buffers can store a maximum of 256 samples, as each 14-bit DAC sample will fill exactly two bytes of memory in RAM (the 2 unused bits will all be zeros padded on the MSB or LSB of the DAC output value, depending on the \bitfield{DACDAL} bit), and concurrent DAC samples from DAC0 and DAC1 occupy the two least significant and two most significant bytes of each 32-bit block, respectively. However, the user can configure the number of samples the DAC will place in each half of the buffer using the \register{DACxHBS} register:

\begin{equation*}
	N_{s,half} = 1 + \register{DACxHBS}
\end{equation*}

To illustrate this in an example, shown in Figure \ref{fig:dac-dma-example}, assume that the half buffer size is 2 samples (\register{DACxHBS} = 1). When the DMA is enabled by setting both \bitfield{DACxDEN} bits to 1, the DAC DMA will immediately begin reading samples from the first half of the double buffer and placing them into both DACs simultaneously, starting with the first samples, $a_0$ and $b_0$. Note that $a_0$ is the sample for DAC0, while $b_0$ is the sample for DAC1. At this time, the data in the second half of the buffer should be updated with new samples so that the DMA can seamlessly transition to them after the first half is empty. Furthermore, the user cannot access any of the data in the first half of the buffer while the DMA is still reading samples from it.

When the DMA finishes reading the second samples, $a_1$ and $b_1$, from the first half of the buffer, it detects that the first half is empty (because of the value of \register{DACxHBS}). It then switches over to reading samples from the second half, simultaneously allowing the CPU access to the first half and setting the \bitfield{DACBEIF} interrupt flag. When this interrupt flag is asserted, the user must update the samples in the first half with new values before the DAC finishes emptying the second half, else the data in the first half will be repeated. The \bitfield{DACBH} bit indicates to the user which half of the buffer the DAC is currently reading data from, and which half is free to be written to. While the user is updating the data in the first half of the buffer, the DMA continues to read new samples from the second half, which is now unable to be read by the CPU. The DMA reads samples $a_2$, $b_2$, $a_3$, and $b_3$ from the second half before switching back to the first half. This process continues indefinitely until the DMA is disabled.

\begin{figure}[htpb]
	\centering
	\begin{subfigure}[b]{\linewidth}
		\begin{bytefield}[endianness=little, bitwidth=0.1\linewidth]{10}
			\bitheader[lsb=0]{0-9} \\
			\bitbox{2}{$a_0$} & \bitbox{2}{$b_0$} & \bitbox{2}{$a_1$} & \bitbox{2}{$b_1$} & \bitbox{2}{Not used...} \\
			\bitbox[t]{1}{\tiny \texttt{0x11800}} & \bitbox[t]{1}{\tiny \texttt{0x11801}} & \bitbox[t]{1}{\tiny \texttt{0x11802}} & \bitbox[t]{1}{\tiny \texttt{0x11803}} & \bitbox[t]{1}{\tiny \texttt{0x11804}} & \bitbox[t]{1}{\tiny \texttt{0x11805}} & \bitbox[t]{1}{\tiny \texttt{0x11806}} & \bitbox[t]{1}{\tiny \texttt{0x11807}} & \bitbox[t]{1}{\tiny \texttt{0x11808}} & \bitbox[t]{1}{\tiny \texttt{0x11809}} \\
		\end{bytefield}
		\caption{First Half Buffer (Ping Buffer)}
		\label{fig:dac-dma-example-first-half-buffer}
	\end{subfigure}
	\vspace*{0.25in}

	\begin{subfigure}[b]{\linewidth}
		\begin{bytefield}[endianness=little, bitwidth=0.1\linewidth]{10}
			\bitheader[lsb=0]{0-9} \\
			\bitbox{2}{$a_2$} & \bitbox{2}{$b_2$} & \bitbox{2}{$a_3$} & \bitbox{2}{$b_3$} & \bitbox{2}{Not used...} \\
			\bitbox[t]{1}{\tiny \texttt{0x11C00}} & \bitbox[t]{1}{\tiny \texttt{0x11C01}} & \bitbox[t]{1}{\tiny \texttt{0x11C02}} & \bitbox[t]{1}{\tiny \texttt{0x11C03}} & \bitbox[t]{1}{\tiny \texttt{0x11C04}} & \bitbox[t]{1}{\tiny \texttt{0x11C05}} & \bitbox[t]{1}{\tiny \texttt{0x11C06}} & \bitbox[t]{1}{\tiny \texttt{0x11C07}} & \bitbox[t]{1}{\tiny \texttt{0x11C08}} & \bitbox[t]{1}{\tiny \texttt{0x11C09}} \\
		\end{bytefield}
		\caption{Second Half Buffer (Pong Buffer)}
		\label{fig:dac-dma-example-second-half-buffer}
	\end{subfigure}
	\caption{DAC DMA Example}
	\label{fig:dac-dma-example}
\end{figure}

Note that the end of first half of the buffer is not necessarily adjacent in memory to the second half. In the above example, $a_0$ and $b_0$ to $a_1$ and $b_1$ are in consecutive memory locations, but then the address of $a_2$ jumps up by 512 bytes. The only case where the last sample of the first half is adjacent to the first sample of the second half is when the half buffer size is at its maximum value of 512.



\subsection{A Note on Shared DMA Memory}

The RAM memory containing the ADC and DAC DMA double buffers can be changed by hardware and, at times, become inaccessible to the CPU when the DAC DMA is activated. In this case, the stack pointer should be initialized to a memory address lower than the beginning of these DMA buffers, else the stack will be corrupted by the ADC or DAC data. However, if the ADC and DAC DMAs are never used, then the user is permitted and encouraged to take advantage of the extra RAM space would otherwise have been taken by the double buffers. This can be achieved by modifying the initial stack pointer location to the beginning of the memory segment of the double buffers.



\subsection{Flags and Interrupts}

The \peripheral{DAC} peripheral includes one indicator bit and one interrupt flag in the DAC status register \register{DACxSR}.

The \bitfield{DACBH} indicator bit shows which half of the double buffer the DAC DMA is currently reading samples from. By extension, it also shows that the other buffer half is available for the CPU to write. This bit should be used to determine which half to place new samples into when the DAC buffer half empty interrupt flag \bitfield{DACBEIF} triggers.

The DAC buffer half empty interrupt flag \bitfield{DACBEIF} is set as soon as the DMA finishes reading samples from the current half of the double buffer and goes on to begin reading samples from the next. \bitfield{DACBH} is updated simultaneously. When the DAC ISR executes, \bitfield{DACBH} indicates which half of the double buffer should be updated with new samples. Write any value to the DAC status register \register{DACxSR} to clear this interrupt flag.
